% GENERATED FILE. Edit canonical lessons, then run npm run generate:books.
% canonical-source-hash: fnv1a64:08955fd0cf6e8679
% canonical-lessons: RU-C213-granitsa, RU-C213-tamozhnya, RU-C213-stolitsa, RU-C213-kvartira, RU-C213-etazh

\chapter{Border, Customs, Capital, Flat, Floor}
\label{ch:ru-border-customs-capital-flat-floor}
\hlchaptermodality{213}

\begin{chapteropening}
\textbf{By the end of this chapter:} \emph{I can say a border, customs, a capital, a flat and a floor.}

Complete the last lesson of chapter 213: I can say a border, customs, a capital, a flat and a floor.
\end{chapteropening}

\section[\texorpdfstring{\ru{граница} (granítsa)}{граница (granítsa)}]{\ru{граница} (granítsa) — a border}

\label{lesson:RU-C213-granitsa}



\begin{quote}
Before the new one: say the Russian for September, then the Russian for October.
\end{quote}

\subsection*{You'll want to know: \ru{граница}}
\textbf{\ru{граница}} (\emph{granítsa}) — \textquotedblleft{}a border\textquotedblright{}.

\begin{cousinweb}
\textbf{\ru{граница}} is built on \emph{\ru{грань}}, \textquotedblleft{}an edge, a facet\textquotedblright{} (Vasmer), which first probably meant \textquotedblleft{}a point\textquotedblright{}; Vasmer relates it to Old High German \emph{grana}, \textquotedblleft{}whisker\textquotedblright{}. German \emph{Grenze}, \textquotedblleft{}border\textquotedblright{}, was borrowed from this Slavic word.
\end{cousinweb}

\begin{grammarlens}[title={What you've built}]
One more word for this chapter.
\end{grammarlens}

\subsection*{Your turn}
\begin{itemize}
\raggedright
  \item \emph{Say it:} \emph{granítsa}
  \item \emph{Say it:} \emph{granítsa}, once more
  \item \emph{Say it:} say \emph{granítsa}
  \item \emph{Recall:} say the Russian for September, then the Russian for October, then say \emph{granítsa} again
\end{itemize}

\begin{tcolorbox}[breakable,colback=teal!4,colframe=teal!35!black,title={Before you move on}]
What does \ru{граница} mean? (\textquotedblleft{}A border\textquotedblright{}.) Say it once more.
\end{tcolorbox}
\section[\texorpdfstring{\ru{таможня} (tamózhnya)}{таможня (tamózhnya)}]{\ru{таможня} (tamózhnya) — customs}

\label{lesson:RU-C213-tamozhnya}



\begin{quote}
Before the new one: say the Russian for October, then the Russian for a border.
\end{quote}

\subsection*{You'll want to know: \ru{таможня}}
\textbf{\ru{таможня}} (\emph{tamózhnya}) — \textquotedblleft{}customs\textquotedblright{}.

\begin{cousinweb}
\textbf{\ru{таможня}} comes from Old Russian \emph{\ru{тамга}}, \textquotedblleft{}a stamp, a customs duty\textquotedblright{} (Vasmer), from Turkic \emph{tamga}, \textquotedblleft{}an owner's brand, a seal, a tax\textquotedblright{}.
\end{cousinweb}

\begin{grammarlens}[title={What you've built}]
One more word for this chapter.
\end{grammarlens}

\subsection*{Your turn}
\begin{itemize}
\raggedright
  \item \emph{Say it:} \emph{tamózhnya}
  \item \emph{Say it:} \emph{tamózhnya}, once more
  \item \emph{Say it:} say \emph{tamózhnya}
  \item \emph{Recall:} say the Russian for October, then the Russian for a border, then say \emph{tamózhnya} again
\end{itemize}

\begin{tcolorbox}[breakable,colback=teal!4,colframe=teal!35!black,title={Before you move on}]
What does \ru{таможня} mean? (\textquotedblleft{}Customs\textquotedblright{}.) Say it once more.
\end{tcolorbox}
\section[\texorpdfstring{\ru{столица} (stolítsa)}{столица (stolítsa)}]{\ru{столица} (stolítsa) — a capital (city)}

\label{lesson:RU-C213-stolitsa}



\begin{quote}
Before the new one: say the Russian for a border, then the Russian for customs.
\end{quote}

\subsection*{You'll want to know: \ru{столица}}
\textbf{\ru{столица}} (\emph{stolítsa}) — \textquotedblleft{}a capital (city)\textquotedblright{}.

\begin{cousinweb}
\textbf{\ru{столица}} is built on \emph{\ru{стол}}, which in Old Russian meant \textquotedblleft{}a throne, a princely seat\textquotedblright{} as well as \textquotedblleft{}a table\textquotedblright{} (Vasmer): the capital is where the ruler sits.
\end{cousinweb}

\begin{grammarlens}[title={What you've built}]
One more word for this chapter.
\end{grammarlens}

\subsection*{Your turn}
\begin{itemize}
\raggedright
  \item \emph{Say it:} \emph{stolítsa}
  \item \emph{Say it:} \emph{stolítsa}, once more
  \item \emph{Say it:} say \emph{stolítsa}
  \item \emph{Recall:} say the Russian for a border, then the Russian for customs, then say \emph{stolítsa} again
\end{itemize}

\begin{tcolorbox}[breakable,colback=teal!4,colframe=teal!35!black,title={Before you move on}]
What does \ru{столица} mean? (\textquotedblleft{}A capital (city)\textquotedblright{}.) Say it once more.
\end{tcolorbox}
\section[\texorpdfstring{\ru{квартира} (kvartíra)}{квартира (kvartíra)}]{\ru{квартира} (kvartíra) — a flat}

\label{lesson:RU-C213-kvartira}



\begin{quote}
Before the new one: say the Russian for customs, then the Russian for a capital.
\end{quote}

\subsection*{You'll want to know: \ru{квартира}}
\textbf{\ru{квартира}} (\emph{kvartíra}) — \textquotedblleft{}a flat\textquotedblright{}.

\begin{cousinweb}
\textbf{\ru{квартира}} — from 1704, also folk \emph{\ru{фатера}}; it came through German \emph{Quartier} or Dutch \emph{kwartier} from Old French \emph{quartier}, \textquotedblleft{}a quarter, lodgings\textquotedblright{}, Latin \emph{quartārius}, \textquotedblleft{}a fourth\textquotedblright{} (Vasmer).
\end{cousinweb}

\begin{grammarlens}[title={What you've built}]
One more word for this chapter.
\end{grammarlens}

\subsection*{Your turn}
\begin{itemize}
\raggedright
  \item \emph{Say it:} \emph{kvartíra}
  \item \emph{Say it:} \emph{kvartíra}, once more
  \item \emph{Say it:} say \emph{kvartíra}
  \item \emph{Recall:} say the Russian for customs, then the Russian for a capital, then say \emph{kvartíra} again
\end{itemize}

\begin{tcolorbox}[breakable,colback=teal!4,colframe=teal!35!black,title={Before you move on}]
What does \ru{квартира} mean? (\textquotedblleft{}A flat\textquotedblright{}.) Say it once more.
\end{tcolorbox}
\section[\texorpdfstring{\ru{этаж} (etázh)}{этаж (etázh)}]{\ru{этаж} (etázh) — a floor, a storey}

\label{lesson:RU-C213-etazh}



\begin{quote}
Before the new one: say the Russian for a capital, then the Russian for a flat.
\end{quote}

\subsection*{You'll want to know: \ru{этаж}}
\textbf{\ru{этаж}} (\emph{etázh}) — \textquotedblleft{}a floor, a storey\textquotedblright{}.

\begin{cousinweb}
\textbf{\ru{этаж}} comes from French \emph{étage}, \textquotedblleft{}a storey\textquotedblright{}, from a Latin \emph{*staticum}, from \emph{statiō}, \textquotedblleft{}a standing, a post\textquotedblright{} (Vasmer).
\end{cousinweb}

\begin{grammarlens}[title={What you've built}]
One more word for this chapter.
\end{grammarlens}

\subsection*{Your turn}
\begin{itemize}
\raggedright
  \item \emph{Say it:} \emph{etázh}
  \item \emph{Say it:} \emph{etázh}, once more
  \item \emph{Say it:} say \emph{etázh}
  \item \emph{Recall:} say the Russian for a capital, then the Russian for a flat, then say \emph{etázh} again
\end{itemize}

\begin{tcolorbox}[breakable,colback=teal!4,colframe=teal!35!black,title={Before you move on}]
What does \ru{этаж} mean? (\textquotedblleft{}A floor, a storey\textquotedblright{}.) Say it once more.
\end{tcolorbox}
